% ==================================================================
\section{How LLMs Change the Value of Failure}
\label{sec:opportunity}
% ==================================================================

The arrival of large language models has changed this calculus in
two directions simultaneously. On one hand, LLMs have lowered the
barrier to producing manuscripts, accelerating the very pressures
described above. The volume of submissions to major conferences
reflects this. ICLR received roughly 7,000 submissions in 2024,
12,000 in 2025, and nearly 20,000 in 2026
\citep{pangram2025iclr}. On the other hand, LLMs have made it
practical for the first time to publish, retrieve, and exploit
failure at scale. This section examines three ways in which this
changes the value of failure data.

\subsection{LLMs Have Relaxed the Retrieval Bottleneck}

A model processing over a hundred thousand words in seconds can
retrieve, compare, and synthesize recorded knowledge at a scale no
human reader ever could. This directly addresses the economic logic
that justified the file drawer. If failures were systematically
published, LLMs could surface them before a researcher commits
months to a direction that others have already explored and
abandoned. The cost of retrieval, which once exceeded the cost of
simply trying again, has dropped by orders of magnitude. The triage
that produced the file drawer is no longer necessary.

At the collective level, this shift compounds. The redundant
failures described in Section~\ref{sec:background}, invisible
because no individual researcher could track what every other group
had tried, become visible when an LLM can cross-reference failure
records across labs, disciplines, and languages simultaneously.
What was an unavoidable collective tax on scientific progress
becomes an avoidable inefficiency.

\subsection{LLMs Need Less Biased Training Data}

From a training perspective, LLMs face a data crisis. Frontier
models have consumed most available high-quality text
\citep{villalobos2024run}, and training on model-generated text
leads to distributional collapse \citep{shumailov2024curse}. The
field needs new, human-generated data. The largest untapped corpus
meeting this criterion is precisely the one that was never
published.

Beyond volume, the composition of training data matters. An LLM
trained predominantly on positive-result literature inherits a
distorted model of science in which most experiments succeed. This
positive skew may systematically miscalibrate the model's judgment
about which research directions are viable. Incorporating published
failure data would not only expand the training corpus but also
correct the distributional bias at its source.

Early evidence from domain-specific applications supports this.
\citet{toniato2025negative} incorporated failed chemical reactions
into language model training and found improved prediction accuracy,
demonstrating that negative reactions carry learnable signal because
each attempt is grounded in background knowledge.
\citet{lee2025banel} went further, showing that a generative model
can be improved using only failed attempts, with no positive
examples at all, achieving orders-of-magnitude gains in success
rate on sparse-reward tasks. \citet{naser2025failure} surveyed this
emerging direction more broadly, documenting techniques such as
negative knowledge distillation and error-based curriculum learning
that are designed to extract signal from suboptimal outcomes. These
results indicate that the methodological infrastructure to exploit
failure data already exists. What is missing is not the ability to
learn from failure but the published failure data to learn from.

\subsection{LLM Reviewers Need Failure to Judge Well}

The peer review system is already straining. At ICLR 2026, 21\% of
75,800 peer reviews were flagged as entirely AI-generated
\citep{pangram2025iclr}. At ICML 2025, papers were found to
contain hidden prompts designed to manipulate LLM reviewers into
giving favorable scores \citep{theocharopoulos2025prompt}. Multiple
major venues have responded with escalating bans and detection
measures in their calls for papers. But if the volume of
submissions continues to grow at the current rate, human-only
review will not scale. The question is not whether LLMs will
participate in peer review, but whether they will do so with the
judgment necessary to catch flawed work.

A reviewer, human or machine, trained exclusively on a literature
in which positive results exceed 85\%
\citep{fanelli2012negative} may lack a well-calibrated model of
what failure looks like. We hypothesize that exposure to structured
failure data would improve the ability of LLM reviewers to detect
methodological errors and implausible claims. This remains
untested, and we propose a direct experiment in
Section~\ref{sec:future}.

Current responses are predominantly defensive, deploying detection
tools in an escalating arms race against increasingly sophisticated
fabrication. A complementary approach is structural. Rather than
attempting to prevent bad science from entering the system, building
legitimate channels through which negative results can be published
and absorbed would equip both human and machine reviewers with a
more accurate model of scientific reality. The goal is not to block
the production of waste but to eliminate the conditions that make
waste the rational output.